\begin{frame}
    \frametitle{Objetivos del libro}  
    \begin{itemize}
        \item Presentar la evolución reciente de las finanzas públicas
        \item Dimensionar la magnitud de la problemática presente y futura 
        \item Proponer soluciones para recuperar la sostenibilidad fiscal
    \end{itemize}
\end{frame}

\begin{frame}
    \frametitle{Mensajes clave}  
\begin{block}{Reformas legales y constitucionales}
La magnitud del ajuste fiscal requerido desborda las capacidades del gobierno bajo el marco legal e institucional vigente, lo que hace indispensables reformas legales y constitucionales.
\end{block}
\begin{block}{El gasto público difícilmente podrá reducirse significativamente.}
Las reformas al gasto son necesarias, pero sus efectos se reflejarán más en la contención del crecimiento del gasto que en recortes significativos.
\end{block}
\begin{block}{Se necesita al menos una reforma tributaria, pero...}
Se requiere de un aumento significativo de los ingresos tributarios, lo cual necesita como prerequisito reflexionar sobre el fracaso de las reformas tributarias recientes.
\end{block}
\end{frame}

\begin{frame}{Metodología y fuentes}
\begin{itemize}
    \item Análisis centrado en el Gobierno Nacional Central (GNC).
    \item Cifras en \% del PIB. 0,1\% equivale a $\sim$\$1,8 billones en 2025.
    \item Datos oficiales ajustados para claridad y consistencia.
    \item Fuentes: MinHacienda, DIAN y CARF, principalmente.
    \item Material disponible en: \url{https://sites.google.com/site/oliverpardo/}
\end{itemize}
\end{frame}
